\documentclass{llncs}
\usepackage{graphicx}
\usepackage{booktabs}

\begin{document}

% Link to the documentation usage https://texdoc.org/serve/llncs/0

\title{Player Tracking in Sports --- Multi-View Tracking and 3D Reconstruction}
\author{Andrea Blushi\inst{1} (ID: 264468) \and Davide Don\`{a}\inst{1} (ID: 264467)}
\institute{University of Trento, Master's Degree in Computer Science, Italy\\}
\maketitle

\begin{abstract}
Reliable player and ball tracking in multi-camera sports footage is a prerequisite
for automated game analysis, yet it remains challenging due to rapid motion,
frequent mutual occlusions, and domain-specific visual characteristics such as
specular floor reflections.
We present an end-to-end pipeline applied to synchronized three-camera basketball
recordings from the Sanbapolis facility in Trento, Italy, covering detection,
multi-object tracking, and 3D reconstruction.
For detection, we systematically progress from classical background subtraction
(MOG2) through a pre-trained YOLOv11 baseline to a domain-fine-tuned YOLOv26 with
multi-pass inference and class-agnostic non-maximum suppression, achieving the best
detection F1 across all views.
For tracking, we compare SORT against an extended DeepSORT with OSNet re-identification,
ByteTrack low-confidence recovery, and trajectory smoothing, demonstrating improved
identity maintenance through player crossings.
3D player and ball positions are reconstructed per frame via SVD-based triangulation
across the three calibrated views and refined with an RTS-smoothed Kalman filter;
residual error is dominated by calibration noise and occlusion-driven gaps rather
than the triangulation method itself.
\end{abstract}

\keywords{multi-object tracking \and player tracking \and 3D reconstruction \and
multi-view geometry \and person re-identification \and YOLO}

\section{Introduction}

Vision-based player tracking offers a scalable alternative to the expensive dedicated hardware traditionally used in sports analytics. 
However, basketball poses severe challenges for these systems \cite{survey}: rapid movements cause frequent mutual occlusions, 
the fast-moving ball suffers from severe motion blur, and indoor environments introduce specular floor reflections. 
Multi-camera setups alleviate occlusions by providing complementary viewpoints, but they require robust cross-view association, 
which in turn depends on consistent player tracking within each individual camera.

In this work, we present an end-to-end pipeline for multi-view player tracking and 3D reconstruction, 
applied to basketball footage from the Sanbapolis facility in Trento, Italy. 
We systematically compare detection strategies: from classical background subtraction (MOG2), to a domain-adapted YOLOv26 
with multi-pass inference—showing that domain adaptation combined with class-agnostic NMS \cite{nms} is essential for reliable detection. 
For tracking, evaluating SORT \cite{sort} against DeepSORT \cite{deepsort} demonstrates that appearance-based re-identification significantly 
reduces identity switches during player crossings. 
Finally, we reconstruct 3D trajectories via SVD-based triangulation across three calibrated cameras, utilizing Kalman smoothing to suppress noise.
\section{Related Work}\label{sec:related}

\subsection{Background Subtraction}
% MOG2 family methods

\subsection{Object Detection with Deep Learning}
% YOLO lineage

\subsection{Multi-Object Tracking}
% SORT, DeepSORT, Re-ID models

\subsection{Multi-View 3D Reconstruction in Sports}
% Prior work on camera-based sports analytics

\section{Dataset and Experimental Setup}\label{sec:dataset}

% Sanbapolis facility dataset — 3 synchronized cameras (cam_13, cam_2, cam_4)
% Video specs: 25 fps, synchronized recordings
% Ground truth: Roboflow annotations, sparse (every 5 frames)
% Key hyperparameters table

\section{Detection}\label{sec:detection}

\subsection{Background Subtraction (MOG2 + CLAHE)}
As a baseline, we adopt MOG2 background subtraction \cite{survey}, in which each
pixel's intensity distribution is modelled as a mixture of $K$ Gaussians:
\begin{equation}
  P(x_t) = \sum_{k=1}^{K} w_{k,t}\,\mathcal{N}(x_t \mid \mu_{k,t},\,\sigma_{k,t}^2).
\end{equation}
A pixel is classified as foreground if its Mahalanobis distance to the
best-matching background component exceeds a fixed threshold.
Prior to background modelling, illumination is normalised with Contrast Limited
Adaptive Histogram Equalization (CLAHE) \cite{clahe}, and the resulting foreground
mask is cleaned with morphological opening and closing; connected components smaller
than a minimum area threshold are discarded.
Alternative background subtractors (KNN, LSBP, GSOC) were evaluated but offered no
improvement over MOG2 while incurring higher computational cost.

Despite these refinements, the method exhibits three recurring failure modes: floor
specularities produce shadows whose hue values fall outside the range that MOG2's
built-in shadow detector handles; players who are briefly stationary are
indistinguishable from background; and moving spectators generate persistent false
positives.
Because the optimal morphological kernel sizes and detection thresholds differ
across cameras, the method also lacks cross-camera generalization.

\subsection{Pre-Trained YOLOv11}
To overcome these limitations, we apply a YOLOv11 model \cite{yolo_paper,yolo_github}
pre-trained on COCO, using a confidence threshold of 0.4, an NMS IoU threshold of
0.45, and an inference resolution of 640\,px.
Detection quality improves substantially, with fewer false positives and better
spatial localisation.
However, the pre-trained model has two domain gaps: the ball is absent from COCO's
person-centric classes, and players and referees are both labelled ``person'',
making identity distinction impossible.

\subsection{Fine-Tuned YOLOv26 with Multi-Pass Inference}
To close both gaps, we fine-tune the YOLOv26 backbone on a Sanbapolis-specific
annotated set.
The fine-tuned model uses identity-encoding class labels (one class per player
identity, e.g.\ \textit{Red\_11}) and a dedicated \textit{Ball} class.

Because at 640\,px the convolutional backbone downsamples a 1080p frame by
approximately $1.7\times$, a 15\,px basketball occupies fewer than two feature-map
grid cells, which is insufficient for reliable localisation.
We therefore run two separate inference passes per frame:
\begin{itemize}
  \item \textbf{Player pass}: 1280\,px resolution, confidence threshold 0.3.
    A second run at a lower threshold (0.15) retains a low-confidence band that is
    reserved for the tracker's secondary association pass.
  \item \textbf{Ball pass}: 1600\,px resolution, confidence threshold 0.2.
    The higher resolution increases the spatial footprint of the ball, while the
    lower threshold favours recall for this motion-blurred, occasionally occluded object.
\end{itemize}
The two sets of detections are merged frame by frame into a single output.

\subsubsection{Class-Agnostic NMS}
Because player identity is encoded directly in the class label, the standard
per-class NMS may retain two overlapping boxes that predict different identities for
the same physical player.
We therefore replace it with a class-agnostic variant \cite{nms}: all detections are
pooled regardless of class, sorted by confidence, and greedily suppressed so that
any box whose Intersection over Union with a higher-scoring box exceeds 0.75 is
discarded.
This collapses duplicate boxes for the same player without affecting the ball.
The result of this step is the final detection output fed to the tracker.

\begin{figure}[t]
  \centering
  \includegraphics[width=0.95\linewidth]{../media/fig_detection.png}
  \caption{Sample detection output from fine-tuned YOLOv26 with class-agnostic NMS
           on cam\_13 (three selected frames). Colour-coded bounding boxes indicate
           player identities; the ball is shown separately.}
  \label{fig:detection}
\end{figure}

\section{Tracking}\label{sec:tracking}

\subsection{SORT}
% Kalman filter state: position, velocity, scale, aspect ratio
% Hungarian matching on IoU; max IoU dist=0.8, max age=60 frames, n_init=2
% Limitation: identity switches during player crossings/occlusions

\subsection{DeepSORT}
% Enhancement: OSNet Re-ID embeddings (cosine distance)
% Matching cascade: recency -> appearance gate -> IoU fallback
% Gallery: up to 100 embeddings per track
% Parameters: max appearance dist=0.2, max age=30 frames, n_init=2
% Key advantage: fewer identity switches via appearance-based reassociation

\subsection{Label Resolution}
% Voting mechanism: assign stable class label per track ID across frames
% Collision resolution when multiple tracks claim same identity
% Ball is exempt (categorical object, not a unique player identity)

\section{3D Reconstruction}\label{sec:reconstruction}

\subsection{Camera Calibration}
% Intrinsics: provided directly
% Extrinsics: computed from user-clicked court landmarks paired with known 3D court coords
% Distortion modeled directly (no pre-undistortion step)

\subsection{Rectification}
% Removes lens distortion from 2D detections
% Maps bounding box centers to a common undistorted coordinate system

\subsection{Triangulation}
% Groups 2D detections by track ID and frame index across all cameras
% Builds per-camera projection matrices P = K[R|t]
% SVD-based homogeneous solution, normalized to Euclidean coordinates
% Output: 3D trajectories for each player and the ball

\subsection{Trajectory Smoothing}
% Per-player Kalman filter on 3D coordinates
% State: position + velocity; mitigates triangulation noise and detection jitter

\section{Evaluation}\label{sec:evaluation}

\subsection{Detection Metrics}
% Precision, Recall, F1, Mean IoU @ IoU=0.5
% Evaluated on annotated frames (every 5 frames)
% Compared: MOG2 vs YOLO pretrained vs YOLO fine-tuned

\subsection{Tracking Metrics}

\subsubsection{IDF1-Based Identity Metrics}
% IDTP, IDFP, IDFN -> IDP, IDR, IDF1
% Measures ability to maintain consistent identities over time

\subsubsection{HOTA}
% DetA: TP / (TP + FP + FN) at multiple IoU thresholds
% AssA: per-pair association consistency across frames
% HOTA: geometric mean of DetA x AssA (balanced, neither dominates)
% LocA: mean IoU of matched pairs
% Evaluated at 10 IoU thresholds (0.05 to 0.95), score averaged uniformly
% Compared: SORT vs DeepSORT

\subsection{3D Geometry Metrics}

\subsubsection{Reprojection Error}
% Project triangulated 3D points back to image plane
% Metrics: mean, median, std, RMSE; accuracy @ 5/10/20 px

\subsubsection{Trajectory Error}
% Compare projected 3D trajectories vs GT 2D annotated trajectories
% ADE: average displacement error over all frames
% FDE: final displacement error (captures drift)
% MTE: median trajectory error (robust to outliers)
% Fragments: coverage gaps in predicted vs GT trajectory

\section{Results and Discussion}\label{sec:results}

\subsection{Detection}
% Fine-tuned YOLO + class-agnostic NMS wins; explain why each baseline falls short
% MOG2: low precision/recall from domain violations
% Pretrained YOLO: good generalization but misses ball and player identity
% Fine-tuned YOLO: best across all metrics; NMS essential to remove duplicates

\subsection{Tracking}
% DeepSORT wins on IDF1 and HOTA; quantify reduction in identity switches
% Discuss sensitivity of max_age parameter for basketball crossings

\subsection{3D Reconstruction}
% Reprojection error results and what drives residual error
% Impact of Kalman smoothing on trajectory metrics (ADE, FDE, MTE)
% Remaining error sources: calibration noise, occlusion-driven gaps

\subsection{Limitations and Failure Cases}
% Cases where identities still switch (prolonged occlusion, cluster crossings)
% Ball detection failures (motion blur, occlusion under players)
% Cameras with partial court coverage

\section{Conclusion}\label{sec:conclusion}

% Summary of the full pipeline and findings
% Best configuration: fine-tuned YOLO -> class-agnostic NMS -> DeepSORT -> triangulation -> Kalman smoothing
% Future work: online/automatic calibration, learned triangulation, extended ReID gallery, real-time deployment


\bibliographystyle{splncs04}
\bibliography{references}

\end{document}
